\begin{proof}[Proof of Proposition~\ref{coates:prop:local_MMLP}]
Let us define the \emph{$P$-height} of a lattice point $p \in P$ to be the non-negative rational number $r$ such that $p$ lies on the boundary of $rP$. The set of $P$-heights of points in $P \cap N$ is a finite subset of $[0,1]$. Fix a set of residual points $\mathcal{R}$ for $P$, and write $\mathcal{R}_{\geq h}$ for the set of points in $\mathcal{R}$ of $P$-height at least $h$. Set $f = \sum_{v \in P \cap N} a_v \bm{x}^v$. We will prove the following statement by descending induction on $h$:
  \begin{equation}
  \tag{$\ddagger$}
  \label{coates:induction statement}
  \begin{minipage}{0.8\textwidth}
    The coefficients $a_v$ such that $v$ has $P$-height $h$ are affine-linear functions of the coefficients $a_w$ for $w \in \mathcal{R}_{\geq h}$, and these coefficients $a_w$ are independent.
  \end{minipage}
  \end{equation}
  \noindent To prove the Proposition, we need to prove statement \eqref{coates:induction statement} for all $h \geq 0$. It holds trivially for $h=0$.

  \subsubsection*{(ii) Base case: $h=1$} Points of $P$-height $1$ lie on the boundary of $P$. Consider an edge $e$ of $P$. Without loss of generality we may assume that $e$ is horizontal and at positive vertical height $r$ above the origin. Up to overall multiplication by a monomial, the coefficients of $f$ along $e$ are
  \begin{equation}
    \label{coates:eq:edge coefficients}
    1 + b_1 x + \cdots + b_{d-1} x^{d-1} + x^d  
  \end{equation}
  where $d = n_e r+s$ and $0 \leq s < r$, for some $b_1,\ldots,b_{d-1}$. The mutability condition gives that 
  \begin{equation}
    \label{coates:eq:edge equation}
    1 + b_1 x + \cdots + b_{d-1} x^{d-1} + x^d = 
    (1+x)^{r n_e} \sum_{i=0}^s c_i x^i
  \end{equation}
  for some coefficients $c_i$. This forces $c_0 = c_s = 1$. If $s = 0$ or $s = 1$ then there are no residual points on $e$, and \eqref{coates:induction statement} holds as \eqref{coates:eq:edge coefficients} coincides with $(1+x)^d$. If $s>1$ then to prove \eqref{coates:induction statement} it suffices to show that we can solve \eqref{coates:eq:edge equation} uniquely for $c_i$ in terms of the coefficients $b_{k_1}, \ldots, b_{k_{s-1}}$ of residual points on $e$. That is, it suffices to show that the $(s-1) \times (s-1)$ matrix with $(i,j)$ entry
  \[
      { n_e r \choose k_i - j}
  \]
  is invertible. But the determinant of this matrix is positive -- indeed it counts certain semi-standard Young tableaux, as can be seen by specializing to $1$ all variables in the Giambelli formula~\textup{([29], eqn A6)} that expresses Schur polynomials in terms of elementary symmetric polynomials. Thus statement \eqref{coates:induction statement} holds for $h=1$. 

  \subsubsection*{(iii) Induction step} Suppose that $H$ satisfies $0 < H < 1$ and that statement \eqref{coates:induction statement} holds for all $h > H$. We will prove \eqref{coates:induction statement} for $h=H$. Let $v \in P$ be a lattice point of $P$-height $H$. Since $P$ is Fano, $v$ lies in the cone $C_e$ over a unique edge $e$. Without loss of generality we may assume that $e$ is horizontal and at positive vertical height $r$ above the origin. Let $d$ denote the lattice length of $e$, and write $d = n_e r + s$ with $0 \leq s < r$. Consider the horizontal line $L$ through $v$; this is at vertical height $Hr$ above the origin. Lattice points on $L$ fall into three classes: residual points in $C_e$, non-residual points in $C_e$, and points outside $C_e$. Convexity implies that the points on $L$ outside $C_e$ are at $P$-height greater than $H$: see Figure~\ref{coates:fig5}. By the induction hypothesis, therefore, coefficients of these points are fixed as affine-linear combinations of coefficients of residual points of height greater than $H$. 
  
  Up to overall multiplication by a monomial, the coefficients of $f$ along $L$ are
  \[
    b_0 + b_1 x + \cdots + b_m x^m  
  \]
  for some $m$ and $b_0,\ldots,b_m$. The mutability condition gives that 
  \begin{equation}
    \label{coates:eq:L equation}
    b_0 + b_1 x + \cdots + b_m x^m = 
    (1+x)^{r H n_e } \sum_{i=0}^s c_i x^i
  \end{equation}
  for some coefficients $c_i$. A primitive $T$-cone with lattice height $r$ contains exactly $h$ lattice points at lattice height $h$, if $0 < h < r$, and so if $L \cap P$ contains no residual points and no points outside $C_e$ then $m<r H n_e$. In this case \eqref{coates:eq:L equation} gives that $b_0 = \cdots = b_m = 0$, and statement \eqref{coates:induction statement} holds. Otherwise \eqref{coates:eq:L equation} holds with $s \geq 0$ and $s+1$ equal to the total number of points in $L \cap P$ that are either residual or lie outside $C_e$. Statement \eqref{coates:induction statement} for $h=H$ follows if we can solve \eqref{coates:eq:L equation} uniquely for $c_i$ in terms of the coefficients $b_{k_0}, \ldots, b_{k_s}$ of points on $L$ that are either residual or lie outside $C_e$. For this, we use the same matrix-invertibility argument as above. This completes the induction step, and the proof of Proposition~\ref{coates:prop:local_MMLP}.
\end{proof}
